\chapter*{Resumen}
\addcontentsline{toc}{chapter}{Resumen}

% Esta tesis doctoral contrasta empíricamente las predicciones de un modelo de agencia multitarea sobre los efectos asimétricos de la descentralización hospitalaria en la eficiencia técnica. El modelo teórico, de tres agentes (pagador público, hospital y médico), predice que la descentralización contractual mejora la eficiencia en cantidad de prestaciones pero tiene efectos ambiguos sobre la intensidad de tratamiento, dependiendo del grado de sustituibilidad entre dimensiones del esfuerzo médico (parámetro $\psi_{12}$).

% El análisis empírico utiliza métodos de frontera estocástica (SFA) con cuatro diseños complementarios y datos de panel del sistema hospitalario español (SIAE y CNH) para el período 2010--2023. Las predicciones contrastadas son: P1 (la descentralización mejora la eficiencia en cantidad), P2 (la descentralización aumenta la ineficiencia en intensidad), P3 (los hospitales privados de mercado son más eficientes en cantidad que los concertados), P4 (el efecto diferencial de la descentralización es mayor en servicios quirúrgicos que en médicos), y P* (las ineficiencias en ambas dimensiones no son independientes).

% Los resultados principales son los siguientes. P1 se confirma robustamente en las siete especificaciones estimadas ($\hat{\alpha}_1 < 0$, significativo al 1\%). P2 no se confirma en su formulación estricta de signo ($\hat{\delta}_1 = -1.477$, signo opuesto al predicho), pero la asimetría entre dimensiones ---la descentralización mejora más la eficiencia en cantidad que en intensidad--- es estadísticamente significativa ($z = -2.09$, $p = 0.037$) y consistente con el modelo bajo sustitución moderada. P3 y P4 se confirman en todos los diseños. P* se confirma con un estadístico de razón de verosimilitudes superior a 10.000.

% La principal contribución es la primera contrastación empírica directa de un modelo de agencia multitarea en el sector hospitalario español, con estimación simultánea de fronteras para cantidad e intensidad y un marco teórico que genera predicciones diferenciadas según la estructura organizativa del hospital.



La organización hospitalaria determina quién diseña los incentivos 
del profesional sanitario, qué dimensiones del \textit{output} son 
verificables por el financiador y cómo se distribuyen el riesgo y 
el control entre los agentes que intervienen en la prestación: es, 
por eso, un objeto situado en la intersección entre el diseño 
institucional, la economía de los incentivos y el Derecho sanitario. 
El ordenamiento español ha configurado un sistema hospitalario que 
permite la coexistencia de formas de gestión con grados muy distintos 
de autonomía directiva, régimen de personal y capacidad contractual. 
Esa configuración arranca del reconocimiento constitucional del derecho 
a la protección de la salud y se articula a través de la Ley General 
de Sanidad, la Ley de Cohesión y Calidad del SNS, la Ley 15/1997 
sobre nuevas formas de gestión y el Estatuto Marco del personal 
estatutario. Esta pluralidad organizativa no es un dato meramente 
descriptivo: condiciona de forma sustantiva los resultados en 
eficiencia técnica que el sistema hospitalario puede alcanzar.

Esta tesis analiza las consecuencias de esa arquitectura institucional 
sobre la eficiencia hospitalaria. Para ello, desarrolla un modelo de 
agencia multitarea con tres agentes ---financiador público, hospital 
y médico--- que producen conjuntamente dos \textit{outputs} asimétricos: 
la cantidad de atención prestada (que depende del esfuerzo conjunto 
del hospital y del médico) y la intensidad diagnóstico-terapéutica 
por episodio (que depende exclusivamente del esfuerzo del médico). 
El modelo genera predicciones diferenciadas según la estructura 
organizativa: la descentralización contractual mejora la eficiencia 
en cantidad pero, cuando las tareas del médico son sustitutas, 
deteriora la eficiencia en intensidad. El tipo de pago y el grado 
de complementariedad entre tareas ---aproximado en la tesis por la 
naturaleza quirúrgica o médica del servicio--- moderan ese efecto 
de forma asimétrica.

El contraste empírico utiliza técnicas de frontera estocástica con 
cuatro diseños complementarios y datos de panel del sistema hospitalario 
español para el período 2010--2023. Los resultados confirman que la 
descentralización reduce significativamente la ineficiencia en cantidad 
(en las siete especificaciones de robustez estimadas), aunque sin la 
inversión de signo predicha en la dimensión de intensidad. No obstante, 
la asimetría entre ambas dimensiones es estadísticamente significativa 
y consistente con el modelo bajo sustitución moderada. La proporción 
de actividad quirúrgica ---proxy de alta complementariedad entre 
tareas--- modera asimétricamente el efecto de la estructura 
organizativa, con signos opuestos entre la frontera quirúrgica y la 
médica. Este constituye el resultado individual más robusto del 
análisis. Finalmente, el test conjunto confirma que los determinantes 
institucionales de la ineficiencia operan de forma cualitativamente 
distinta sobre las dos dimensiones del \textit{output} hospitalario.

La contribución de la tesis es triple. En el plano institucional, 
sistematiza el marco jurídico que hace inteligible la diversidad 
organizativa hospitalaria española y muestra que las categorías del 
Derecho sanitario ---forma de gestión, régimen del personal, mecanismo 
de financiación--- son variables analíticamente relevantes para un 
problema de incentivos. En el plano teórico, extiende los modelos 
clásicos de multitarea al caso de tres agentes con dos \textit{outputs} 
asimétricos en un contexto hospitalario. En el plano empírico, 
proporciona la primera contrastación directa de un modelo de agencia 
multitarea en el sector hospitalario español, con estimación simultánea 
de fronteras para cantidad e intensidad que permite captar los efectos 
asimétricos que un modelo de eficiencia multiproducto convencional 
no podría identificar.

Los resultados tienen implicaciones directas para el diseño 
institucional del sistema hospitalario. La evidencia confirma que 
la estructura organizativa produce efectos estadísticamente 
significativos y cualitativamente distintos sobre las dos dimensiones 
del \textit{output}: un enfoque de política de gestión hospitalaria 
que ignore esa multidimensionalidad ---ya sea apostando por la 
centralización total o por la descentralización generalizada--- agrega 
efectos de signo contrario y puede inducir conclusiones erróneas sobre 
el desempeño real de los hospitales. Una política informada por estos 
resultados debería incorporar indicadores separados de eficiencia en 
volumen e intensidad en los contratos-programa autonómicos y en los 
pliegos de los contratos de concesión, y extender el pago por actividad 
ajustado por complejidad a toda la red concertada. Los resultados no 
establecen causalidad de forma concluyente puesto que la asignación de 
hospitales a formas descentralizadas no es aleatoria, pero son 
consistentes con los mecanismos predichos por el modelo teórico y 
robustos frente a las especificaciones alternativas estimadas. En cualquier caso, la tesis constituye un primer paso en la dirección de una agenda de investigación que integre el análisis institucional, teórico y empírico para comprender el desempeño del sistema hospitalario español y orientar su diseño futuro.

\medskip
\noindent\textbf{Palabras clave:} eficiencia hospitalaria, frontera estocástica, teoría de agencia, multitarea, estructura organizativa, sistema hospitalario español, formas de gestión sanitaria, información asimétrica

\medskip
\noindent Clasificación JEL: I11, I18, D82, D86, C23

\chapter*{Abstract}
\addcontentsline{toc}{chapter}{Abstract}

% This doctoral thesis empirically tests the predictions of a multitask agency model on the asymmetric effects of hospital decentralization on technical efficiency. The theoretical model, with three agents (public payer, hospital, and physician), predicts that contractual decentralization improves efficiency in the quantity of services but has ambiguous effects on treatment intensity, depending on the degree of substitutability between dimensions of physician effort (parameter $\psi_{12}$).

% The empirical analysis uses stochastic frontier analysis (SFA) with four complementary designs and panel data from the Spanish hospital system (SIAE and CNH) for the period 2010--2023. The hypotheses tested are: P1 (decentralization improves quantity efficiency), P2 (decentralization increases inefficiency in intensity), P3 (private market hospitals are more quantity-efficient than contracted hospitals), P4 (the differential effect of decentralization is greater in surgical than in medical services), and P* (inefficiencies in both dimensions are not independent).

% The main results are as follows. P1 is robustly confirmed in all seven estimated specifications ($\hat{\alpha}_1 < 0$, significant at 1\%). P2 is not confirmed in its strict sign formulation ($\hat{\delta}_1 = -1.477$, sign opposite to predicted), but the predicted asymmetry between dimensions ---decentralization improves quantity efficiency more than intensity efficiency--- is statistically significant ($z = -2.09$, $p = 0.037$) and consistent with the model under moderate substitution. P3 and P4 are confirmed across all designs. P* is confirmed with a likelihood ratio statistic exceeding 10,000.

% The main contribution is the first direct empirical test of a multitask agency model in the Spanish hospital sector, with simultaneous estimation of frontiers for quantity and intensity and a theoretical framework that generates differentiated predictions according to the organizational structure of the hospital.

Hospital organisation determines who designs the incentives of 
healthcare professionals, which dimensions of \textit{output} are 
verifiable by the payer, and how risk and control are distributed 
among the agents involved in care delivery: it is, for that reason, 
an object situated at the intersection of institutional design, 
the economics of incentives, and health law. The Spanish legal 
framework has configured a hospital system that allows the coexistence 
of management forms with very different degrees of managerial autonomy, 
employment regime, and contractual capacity. That configuration 
originates in the constitutional recognition of the right to health 
protection and is articulated through the General Health Act, the 
Act on Cohesion and Quality of the National Health System, Act 
15/1997 on new management forms, and the Framework Statute for 
statutory healthcare personnel. This organisational plurality is 
not merely descriptive: it substantively conditions the technical 
efficiency outcomes that the hospital system can achieve.

This thesis analyses the consequences of that institutional 
architecture for hospital efficiency. To do so, it develops a 
multitask agency model with three agents ---public payer, hospital, 
and physician--- who jointly produce two asymmetric \textit{outputs}: 
the quantity of care delivered (which depends on the joint effort 
of the hospital and the physician) and the diagnostic-therapeutic 
intensity per episode (which depends exclusively on the physician's 
effort). The model generates differentiated predictions according 
to organisational structure: contractual decentralisation improves 
efficiency in quantity but, when the physician's tasks are 
substitutes, deteriorates efficiency in intensity. The payment 
mechanism and the degree of complementarity between tasks ---proxied 
in the thesis by the surgical or medical nature of the service--- 
moderate that effect asymmetrically.

The empirical test uses stochastic frontier techniques with four 
complementary designs and panel data from the Spanish hospital 
system for the period 2010--2023. The results confirm that 
decentralisation significantly reduces inefficiency in quantity 
(across all seven robustness specifications estimated), albeit 
without the sign reversal predicted for the intensity dimension. 
Nevertheless, the asymmetry between both dimensions is statistically 
significant and consistent with the model under moderate substitution. 
The proportion of surgical activity ---a proxy for high 
complementarity between tasks--- asymmetrically moderates the 
effect of organisational structure, with opposite signs across 
the surgical and medical frontiers. This is the single most robust 
result of the analysis. Finally, the joint test confirms that the 
institutional determinants of inefficiency operate in qualitatively 
distinct ways across the two dimensions of hospital \textit{output}.

The contribution of the thesis is threefold. At the institutional 
level, it systematises the legal framework that makes the 
organisational diversity of Spanish hospitals intelligible and 
shows that the categories of health law ---management form, 
employment regime, financing mechanism--- are analytically relevant 
variables for an incentive problem. At the theoretical level, it 
extends classic multitask models to the case of three agents with 
two asymmetric \textit{outputs} in a hospital context. At the 
empirical level, it provides the first direct test of a multitask 
agency model in the Spanish hospital sector, with simultaneous 
estimation of frontiers for quantity and intensity that allows 
capturing the asymmetric effects that a conventional multi-output 
efficiency model could not identify.

The results have direct implications for the institutional design 
of the hospital system. The evidence confirms that organisational 
structure produces statistically significant and qualitatively 
distinct effects on the two dimensions of \textit{output}: a 
hospital management policy that ignores this multidimensionality 
---whether by advocating full centralisation or wholesale 
decentralisation--- aggregates effects of opposite sign and may 
lead to erroneous conclusions about the actual performance of 
hospitals. A policy informed by these results should incorporate 
separate indicators of efficiency in volume and intensity into 
regional performance contracts and into the terms of concession 
contracts, and extend activity-based payment adjusted for complexity 
to the entire contracted network. The results do not establish 
causality conclusively, since the assignment of hospitals to 
decentralised forms is not random, but they are consistent with 
the mechanisms predicted by the theoretical model and robust across 
the alternative specifications estimated. In any case, the thesis 
constitutes a first step towards a research agenda that integrates 
institutional, theoretical, and empirical analysis to understand 
the performance of the Spanish hospital system and guide its 
future design.

\medskip
\noindent
\textbf{Keywords:} hospital efficiency, stochastic frontier analysis, agency theory, multitasking, organizational structure, Spanish hospital system, health care management formulas, asymmetric information

\medskip
\noindent JEL Classification: I11, I18, D82, D86, C23

\chapter*{Agradecimientos}
\addcontentsline{toc}{chapter}{Agradecimientos}
% [Pendiente de redacción]
\vspace{2cm}
\begin{center}
\textit{[Pendiente de redacción]}
\end{center}

\tableofcontents
\listoftables
\listoffigures
